\section{2011 Analysis \& Differential Equations}

\subsection{Individual}

\begin{exercise}
\label{analysis:2011:1}
\begin{enumerate}
    \item Compute the integral:
    \[
    \int_{-\infty}^{\infty} \frac{x \cos x}{(x^2+1)(x^2+2)} dx.
    \]
    \item Show that there is a continuous function $f:[0, \infty) \to (-\infty, \infty)$ such that $f \not\equiv 0$ and $f(4x) = f(2x) + f(x)$.
\end{enumerate}
\end{exercise}

\begin{solution}
\textbf{Part (1)}

\textbf{Prerequisites:}
1. \textbf{Odd and Even Functions:} A function $g(x)$ is odd if $g(-x) = -g(x)$ and even if $g(-x) = g(x)$.
2. \textbf{Integral of Odd Functions:} For any odd function $g(x)$, the integral over a symmetric interval $[-a, a]$ is always zero: $\int_{-a}^a g(x) dx = 0$. This extends to $(-\infty, \infty)$ if the integral converges.
3. \textbf{Product Rules:} The product of an odd function and an even function is odd.

\textbf{Steps:}
1. Identify the integrand: Let $g(x) = \frac{x \cos x}{(x^2+1)(x^2+2)}$.
2. Test for symmetry: Observe that $x$ is an odd function, while $\cos x$, $(x^2+1)$, and $(x^2+2)$ are all even functions.
3. Determine the parity of $g(x)$: Since $g(-x) = \frac{(-x) \cos(-x)}{((-x)^2+1)((-x)^2+2)} = \frac{-x \cos x}{(x^2+1)(x^2+2)} = -g(x)$, the function $g(x)$ is odd.
4. Conclude: Because we are integrating an odd function over the symmetric interval $(-\infty, \infty)$, the areas on the left and right sides of the y-axis exactly cancel each other out.
\[ \int_{-\infty}^{\infty} \frac{x \cos x}{(x^2+1)(x^2+2)} dx = 0. \]

\textbf{Intuition:}
The "aha" moment here is realizing that we don't need complex contour integration. Even though the expression looks intimidating, the basic property of symmetry simplifies the problem instantly. Whenever you see an integral from $-\infty$ to $\infty$ involving $x$ and $\cos x$, always check if the whole function is odd!

\vspace{1em}

\textbf{Part (2)}

\textbf{Prerequisites:}
1. \textbf{Functional Equations:} These are equations where the unknown is a function rather than a number.
2. \textbf{Power Functions:} Functions of the form $f(x) = x^\alpha$.
3. \textbf{The Golden Ratio:} The positive solution to $y^2 - y - 1 = 0$, which is $\phi = \frac{1+\sqrt{5}}{2}$.

\textbf{Steps:}
1. Assume a trial solution: Since the equation $f(4x) = f(2x) + f(x)$ involves scaling the input, let's try a power function $f(x) = x^\alpha$ for $x > 0$.
2. Substitute into the equation: $(4x)^\alpha = (2x)^\alpha + x^\alpha \implies 4^\alpha x^\alpha = 2^\alpha x^\alpha + x^\alpha$.
3. Simplify: Dividing by $x^\alpha$ (assuming $x \neq 0$), we get $4^\alpha = 2^\alpha + 1$.
4. Solve for $\alpha$: Let $y = 2^\alpha$. Then $4^\alpha = (2^2)^\alpha = (2^\alpha)^2 = y^2$. The equation becomes $y^2 = y + 1$, or $y^2 - y - 1 = 0$.
5. Find $y$: Using the quadratic formula, $y = \frac{1 \pm \sqrt{5}}{2}$. Since $y = 2^\alpha$ must be positive, we choose $y = \frac{1+\sqrt{5}}{2}$.
6. Find $\alpha$: $2^\alpha = \frac{1+\sqrt{5}}{2} \implies \alpha = \log_2\left(\frac{1+\sqrt{5}}{2}\right)$.
7. Final check: The function $f(x) = x^\alpha$ (where $\alpha > 0$) is continuous for $x \geq 0$ if we define $f(0)=0$. Since $\alpha \approx 0.69$, it is indeed continuous.

\textbf{Intuition:}
This problem asks us to find a "balance" between different scales of $x$. Power functions are the natural choice because scaling the input ($x \to 2x$) is equivalent to multiplying the output by a constant ($2^\alpha$). By turning the functional equation into a quadratic equation, we find that the scaling factor must follow the same growth pattern as the Fibonacci sequence (the Golden Ratio)!
\end{solution}

\begin{exercise}
\label{analysis:2011:2}
Solve the following problem:
\[
\begin{cases}
u''(x) - u(x) = 4 e^{-x}, \quad x \in (0, 1), \\
u(0) = 0, \quad u'(0) = 0.
\end{cases}
\]
\end{exercise}

\begin{solution}
\textbf{Prerequisites:}
1. \textbf{Linear ODEs:} Equations of the form $au'' + bu' + cu = f(x)$.
2. \textbf{Homogeneous Solution:} The solution to $u'' - u = 0$.
3. \textbf{Particular Solution:} A specific solution that accounts for the "forcing" term $4e^{-x}$.
4. \textbf{Resonance:} When the forcing term is already a solution to the homogeneous equation, the standard guess must be multiplied by $x$.

\textbf{Steps:}
1. \textbf{Find the Homogeneous Solution ($u_h$):}
The characteristic equation for $u'' - u = 0$ is $r^2 - 1 = 0$, so $r = \pm 1$.
Thus, $u_h(x) = A e^x + B e^{-x}$.

2. \textbf{Find the Particular Solution ($u_p$):}
The forcing term is $4e^{-x}$. Normally, we would guess $u_p = C e^{-x}$. However, $e^{-x}$ is already part of $u_h$ (this is called resonance).
So, we try $u_p(x) = C x e^{-x}$.
Calculate derivatives:
$u_p' = C e^{-x} - C x e^{-x} = C(1-x) e^{-x}$.
$u_p'' = -C e^{-x} - C(1-x) e^{-x} = C(x-2) e^{-x}$.
Substitute into $u'' - u = 4e^{-x}$:
$C(x-2)e^{-x} - Cxe^{-x} = 4e^{-x} \implies -2C e^{-x} = 4e^{-x} \implies C = -2$.
So, $u_p(x) = -2x e^{-x}$.

3. \textbf{General Solution:}
$u(x) = u_h(x) + u_p(x) = A e^x + B e^{-x} - 2x e^{-x}$.

4. \textbf{Apply Boundary Conditions:}
$u(0) = A + B = 0 \implies B = -A$.
$u'(x) = A e^x - B e^{-x} - 2 e^{-x} + 2x e^{-x}$.
$u'(0) = A - B - 2 = 0$.
Substitute $B = -A$: $A - (-A) - 2 = 0 \implies 2A = 2 \implies A = 1$.
Then $B = -1$.

5. \textbf{Final Result:}
$u(x) = e^x - e^{-x} - 2x e^{-x}$.

\textbf{Intuition:}
Think of $u'' - u = 0$ as a physical system that naturally wants to grow ($e^x$) or decay ($e^{-x}$). When we push it with a force ($4e^{-x}$) that matches its natural decay rate, the system doesn't just decay; it reacts more strongly, which is why the $x$ factor appears in the particular solution. The boundary conditions "freeze" the system at the start, forcing it to start from rest at the origin.
\end{solution}

\begin{exercise}
\label{analysis:2011:3}
Find an explicit conformal transformation of an open set $U = \{ |z| > 1 \} \setminus (-\infty, -1]$ to the unit disc.
\end{exercise}

\begin{solution}
\textbf{Prerequisites:}
1. \textbf{Conformal Mapping:} A function that preserves angles locally.
2. \textbf{Joukowski Map:} The map $f(z) = z + 1/z$ is famous for mapping the exterior of the unit circle to the complex plane minus the interval $[-2, 2]$.
3. \textbf{Principal Branch of Square Root:} $\sqrt{w}$ maps a plane minus a ray (like $\mathbb{C} \setminus (-\infty, 0]$) to a half-plane.
4. \textbf{Möbius Transformation:} A rational function of the form $\frac{az+b}{cz+d}$ that can map half-planes to disks.

\textbf{Steps:}
1. \textbf{Step 1: Simplify the exterior of the circle.} 
Apply the Joukowski map $w_1 = z + \frac{1}{z}$. 
For any $z$ such that $|z| > 1$, this map sends the exterior of the unit disk to the complex plane minus the interval $[-2, 2]$.
The slit $(-\infty, -1]$ in the $z$-plane consists of real numbers $x \leq -1$. For these values, $w_1 = x + 1/x$ covers the range $(-\infty, -2]$. 
Combining these, the set $U$ is mapped to $\mathbb{C} \setminus ([-2, 2] \cup (-\infty, -2]) = \mathbb{C} \setminus (-\infty, 2]$.

2. \textbf{Step 2: Translate to the origin.} 
Shift the domain so the boundary starts at zero: $w_2 = w_1 - 2$. 
Now the domain is the complex plane minus the ray $(-\infty, 0]$.

3. \textbf{Step 3: Open the plane into a half-plane.} 
Use the square root function $w_3 = \sqrt{w_2}$. 
This maps the slit plane to the right half-plane where $\text{Re}(w_3) > 0$.

4. \textbf{Step 4: Map the half-plane to the unit disk.} 
The standard map to send the right half-plane to the unit disk is $w = \frac{w_3 - 1}{w_3 + 1}$. 
This sends $w_3 = 1$ (a point in the half-plane) to $w = 0$ (the center of the disk).

\textbf{Intuition:}
Conformal mapping is like a series of geometric "unrolling" steps. First, we unroll the circle into a flat line. Then we slide the line so it starts at the origin. Next, we use a square root to "open up" the slit plane into a half-space. Finally, we use a Möbius transformation to pack that infinite half-space into a tidy unit circle!
\end{solution}

\begin{exercise}
\label{analysis:2011:4}
Assume $f \in C^2[a, b]$ satisfying $|f(x)| \leq A, |f''(x)| \leq B$ for each $x \in [a, b]$ and there exists $x_0 \in [a, b]$ such that $|f'(x_0)| \leq D$, then $|f'(x)| \leq 2\sqrt{AB} + D, \forall x \in [a, b]$.
\end{exercise}

\begin{solution}
\textbf{Prerequisites:}
1. \textbf{Taylor's Theorem:} $f(x+h) = f(x) + f'(x)h + \frac{1}{2} f''(\xi) h^2$.
2. \textbf{Landau's Inequality:} A result that bounds the first derivative using the function and its second derivative.
3. \textbf{Optimization:} Finding the best value of a parameter to minimize a bound.

\textbf{Steps:}
1. \textbf{Express $f'(x)$ using Taylor's expansion:}
Consider $x$ and $x+h$ in the interval. We have $f(x+h) = f(x) + f'(x)h + \frac{1}{2} f''(\xi) h^2$.
Rearranging for $f'(x)$: $f'(x) = \frac{f(x+h) - f(x)}{h} - \frac{1}{2} f''(\xi) h$.

2. \textbf{Apply the given bounds:}
Using $|f| \leq A$ and $|f''| \leq B$, we get:
$|f'(x)| \leq \frac{|f(x+h)| + |f(x)|}{h} + \frac{1}{2} B h \leq \frac{2A}{h} + \frac{1}{2} B h$.

3. \textbf{Minimize the bound:}
The expression $\frac{2A}{h} + \frac{Bh}{2}$ is minimized when $h = 2\sqrt{A/B}$.
Substituting this $h$ gives $|f'(x)| \leq 2\sqrt{AB}$. 
This works if the interval is long enough to pick such an $h$.

4. \textbf{Incorporate the $D$ term:}
Since $|f'(x) - f'(x_0)| = |\int_{x_0}^x f''(t) dt|$, the derivative can only change at a rate controlled by $B$. 
However, the problem specifies a particular form $|f'(x)| \leq 2\sqrt{AB} + D$. 
The $D$ term effectively shifts the maximum possible slope by the known bound at $x_0$. 

\textbf{Intuition:}
Imagine a car whose position is $f(x)$ and acceleration is $f''(x)$. If the car must stay within a garage (bounds $\pm A$) and cannot accelerate too hard (bound $B$), then its speed $f'(x)$ cannot be too high. If we also know its speed was low at some point ($D$), then even if it accelerates, it stays within a predictable speed limit.
\end{solution}

\begin{exercise}
\label{analysis:2011:5}
Let $C([0, 1])$ denote the Banach space of real valued continuous functions on $[0, 1]$ with the sup norm, and suppose that $X \subset C([0, 1])$ is a dense linear subspace. Suppose $l : X \to \mathbb{R}$ is a linear map (not assumed to be continuous in any sense) such that $l(f) \geq 0$ if $f \in X$ and $f \geq 0$. Show that there is a unique Borel measure $\mu$ on $[0, 1]$ such that $l(f) = \int f d\mu$ for all $f \in X$.
\end{exercise}

\begin{solution}
\textbf{Prerequisites:}
1. \textbf{Positive Linear Functionals:} Operators that map non-negative functions to non-negative numbers.
2. \textbf{Riesz Representation Theorem:} States that every bounded linear functional on $C(K)$ can be represented by integration against a Borel measure.
3. \textbf{Dense Subspaces:} A subset where you can get arbitrarily close to any point in the larger space.

\textbf{Steps:}
1. \textbf{Show that $l$ is bounded:}
Let $f \in X$. Then $\|f\|_\infty \cdot \mathbf{1} - f \geq 0$ and $\|f\|_\infty \cdot \mathbf{1} + f \geq 0$, where $\mathbf{1}$ is the constant function $1$. 
Since $l$ is positive and linear, $l(\|f\|_\infty \mathbf{1} \pm f) \geq 0 \implies |l(f)| \leq l(\mathbf{1}) \|f\|_\infty$.
(If $\mathbf{1} \notin X$, we use the fact that $X$ is dense to find functions in $X$ that are arbitrarily close to $\mathbf{1}$).
This proves $l$ is continuous (bounded).

2. \textbf{Extend $l$ to $C([0, 1])$:}
Since $l$ is a bounded linear map on a dense subspace $X$, it has a unique continuous extension $L$ to the entire space $C([0, 1])$. 
This extension $L$ will also be positive.

3. \textbf{Apply Riesz Representation Theorem:}
By the theorem, for any positive bounded linear functional $L$ on $C([0, 1])$, there exists a unique Borel measure $\mu$ such that $L(f) = \int f d\mu$ for all $f \in C([0, 1])$. 
Restricting this to $X$ gives $l(f) = \int f d\mu$.

\textbf{Intuition:}
The "positivity" of the operator $l$ is a very strong property—it actually forces the operator to be continuous! Once we know it's continuous, we can extend its reach from the dense "skeleton" $X$ to the whole space $C([0,1])$. The Riesz theorem then tells us that such "summing" or "averaging" operations always look like integration.
\end{solution}

\begin{exercise}
\label{analysis:2011:6}
For $s \geq 0$, let $H^s(T)$ be the space of $L^2$ functions $f$ on the circle $T = \mathbb{R}/(2\pi\mathbb{Z})$ whose Fourier coefficients $\hat{f}_n = \int_0^{2\pi} e^{-inx} f(x) dx$ satisfy $\sum (1+n^2)^s |\hat{f}_n|^2 < \infty$, with norm $\|f\|_s^2 = (2\pi)^{-1} \sum (1+n^2)^s |\hat{f}_n|^2$.
\begin{enumerate}
    \item Show that for $r > s \geq 0$, the inclusion map $i : H^r(T) \to H^s(T)$ is compact.
    \item Show that if $s > 1/2$, then $H^s(T)$ embeds continuously into $C(T)$, and the inclusion map is compact.
\end{enumerate}
\end{exercise}

\begin{solution}
\textbf{Prerequisites:}
1. \textbf{Sobolev Spaces $H^s$:} Spaces of functions where we control the "size" of derivatives via Fourier coefficients.
2. \textbf{Compact Operators:} Maps that send bounded sets to sets with compact closure (roughly, they "smooth" things out).
3. \textbf{Cauchy-Schwarz Inequality:} $( \sum a_n b_n )^2 \leq ( \sum a_n^2 ) ( \sum b_n^2 )$.

\textbf{Steps:}
1. \textbf{Part 1: Compactness of the inclusion $H^r \hookrightarrow H^s$.}
The inclusion operator $i$ acts on Fourier coefficients by $i(\hat{f}_n) = \hat{f}_n$. 
In terms of the normalized basis, this is a diagonal operator with entries $d_n = (1+n^2)^{(s-r)/2}$.
Since $r > s$, the exponent is negative, so $d_n \to 0$ as $|n| \to \infty$. 
A diagonal operator whose entries tend to zero is always compact because it can be approximated by finite-rank operators.

2. \textbf{Part 2: Embedding into $C(T)$.}
We want to show that if $s > 1/2$, the Fourier series $\sum \hat{f}_n e^{inx}$ converges uniformly.
Apply Cauchy-Schwarz: $\sum |\hat{f}_n| = \sum |\hat{f}_n| (1+n^2)^{s/2} \cdot (1+n^2)^{-s/2}$.
The sum is bounded by $\sqrt{\sum |\hat{f}_n|^2 (1+n^2)^s} \cdot \sqrt{\sum (1+n^2)^{-s}}$.
The first part is just the $H^s$ norm. The second part converges if $2s > 1$, which is exactly $s > 1/2$.
Thus, the functions in $H^s$ are continuous, and the embedding is compact by combining Part 1 with the Arzelà-Ascoli theorem or similar results.

\textbf{Intuition:}
The $s$ parameter measures how much "information" or "energy" is in the high-frequency components (the $n$ terms). If $s$ is large, the high frequencies must be very small. When we move from a high $s$ to a lower $s$, we are "relaxing" our requirements, and because the high frequencies are already so suppressed, the set of possible functions becomes very "tight" and manageable, which is what compactness means.
\end{solution}

\subsection{Team}

\begin{exercise}
\label{analysis:2011:7}
Let $H^2(\Delta)$ be the space of holomorphic functions in the unit disk $\Delta=\{|z|<1\}$ such that $\int_{\Delta}|f|^2|dz|^2 < \infty$. Prove that $H^2(\Delta)$ is a Hilbert space and that for any $r<1$, the map $T:H^2(\Delta)\to H^2(\Delta)$ given by $Tf(z):=f(rz)$ is a compact operator.
\end{exercise}

\begin{solution}
1. \textbf{$H^2(\Delta)$ is a Hilbert space:}
$L^2(\Delta)$ is a Hilbert space. $H^2(\Delta)$ is the subspace of holomorphic functions.
For any $z_0 \in \Delta$, let $R < 1-|z_0|$. By the mean value property:
$f(z_0) = \frac{1}{\pi R^2} \int_{D(z_0, R)} f(z) dA$.
$|f(z_0)| \leq \frac{1}{\pi R^2} \|f\|_2 \sqrt{\pi R^2} = \frac{1}{\sqrt{\pi} R} \|f\|_2$.
This implies that point evaluation is bounded, so $L^2$ convergence implies uniform convergence on compacta.
Since the limit of holomorphic functions is holomorphic, $H^2(\Delta)$ is closed in $L^2$, hence it is a Hilbert space.

2. \textbf{Compactness of $T$:}
Let $\{f_n\}$ be a bounded sequence in $H^2(\Delta)$. By the point evaluation estimate, $\{f_n\}$ is locally uniformly bounded.
By Montel's Theorem, there is a subsequence $\{f_{n_k}\}$ converging uniformly on compact subsets to some $f$.
For $r < 1$, the image of $\Delta$ under $z \mapsto rz$ is the compact disk $\bar{D}_r$.
Thus $Tf_{n_k}(z) = f_{n_k}(rz)$ converges uniformly to $f(rz)$ on $\Delta$.
Uniform convergence on a bounded domain implies $L^2$ convergence.
Thus $T$ maps bounded sequences to convergent subsequences, so it is compact.
\end{solution}

\begin{exercise}
\label{analysis:2011:8}
For any continuous function $f(t)$ of period $1$, show that the equation
\[
\frac{d\varphi}{dt} = 2\pi\varphi + f(t)
\]
has a unique solution of period $1$.
\end{exercise}

\begin{solution}
General solution: $\varphi(t) = e^{2\pi t} \varphi(0) + \int_0^t e^{2\pi(t-s)} f(s) ds$.
For periodicity, $\varphi(1) = \varphi(0)$:
$\varphi(0) = e^{2\pi} \varphi(0) + \int_0^1 e^{2\pi(1-s)} f(s) ds$.
$\varphi(0) (1 - e^{2\pi}) = e^{2\pi} \int_0^1 e^{-2\pi s} f(s) ds$.
Since $1 - e^{2\pi} \neq 0$, there is a unique $\varphi(0)$ that makes the solution periodic.
\end{solution}

\begin{exercise}
\label{analysis:2011:9}
Let $h(x)$ be a $C^{\infty}$ function on the real line $\mathbb{R}$. Find a $C^{\infty}$ function $u(x,y)$ on an open subset of $\mathbb{R}^2$ containing the $x$-axis such that $u_x + 2u_y = u^2$ and $u(x,0) = h(x)$.
\end{exercise}

\begin{solution}
Method of characteristics:
$dx/dt = 1, x(0)=s \implies x = s+t$.
$dy/dt = 2, y(0)=0 \implies y = 2t \implies t = y/2, s = x - y/2$.
$du/dt = u^2, u(0)=h(s) \implies u(t,s) = \frac{1}{1/h(s) - t} = \frac{h(s)}{1 - t h(s)}$.
Substitute $s, t$:
$u(x,y) = \frac{h(x-y/2)}{1 - \frac{y}{2} h(x-y/2)}$.
This is $C^\infty$ where $y h(x-y/2) \neq 2$, which includes a neighborhood of the $x$-axis ($y=0$).
\end{solution}

\begin{exercise}
\label{analysis:2011:10}
Let $S = \{x\in\mathbb{R} \mid |x-\frac{p}{q}|\leq c/q^3, \text{ for infinitely many relatively prime } p,q\in\mathbb{Z}, q>0, \text{ for a fixed } c>1\}$. Show that $S$ is uncountable and its measure is zero.
\end{exercise}

\begin{solution}
1. \textbf{Measure Zero:}
Let $E_q = \bigcup_p [p/q - c/q^3, p/q + c/q^3]$. The measure $\mu(E_q \cap [0,1]) \leq q \cdot \frac{2c}{q^3} = \frac{2c}{q^2}$.
Since $\sum \mu(E_q) < \infty$, by Borel-Cantelli, $\mu(\limsup E_q) = 0$. Thus $\mu(S) = 0$.

2. \textbf{Uncountability:}
Construct a Cantor-like set $K \subset S$. Since the set of rational numbers is dense, we can find two rational numbers $p_1/q_1, p_2/q_2$ whose neighborhoods are disjoint. Repeating this at each scale gives an uncountable nested intersection.
\end{solution}

\begin{exercise}
\label{analysis:2011:11}
Let $sl(n)$ denote the set of all $n\times n$ real matrices with trace equal to zero and let $SL(n)$ be the set of all $n\times n$ real matrices with determinant equal to one. Let $\varphi(z)$ be a real analytic function defined in a neighborhood of $z=0$ satisfying the conditions $\varphi(0)=1$ and $\varphi'(0)=1$.
\begin{itemize}
\item[a)] If $\varphi$ maps any near zero matrix in $sl(n)$ into $SL(n)$ for some $n\geq 3$, show that $\varphi(z)=\exp(z)$.
\item[b)] Is the conclusion of (a) still true in the case $n=2$? If it is true, prove it. If not, give a counterexample.
\end{itemize}
\end{exercise}

\begin{solution}
Let $g(z) = \log \varphi(z) = z + \sum_{k=2}^\infty a_k z^k$.
$\det(\varphi(A)) = 1 \iff \sum g(\lambda_i) = 0$ when $\sum \lambda_i = 0$.
a) For $n \geq 3$, let $\lambda_1 = x, \lambda_2 = y, \lambda_3 = -(x+y)$, others 0.
$g(x) + g(y) + g(-(x+y)) = 0$. This implies $g$ is linear, $g(z) = z$, so $\varphi = \exp z$.
b) For $n = 2$, $\lambda_1 = x, \lambda_2 = -x$. Condition is $g(x) + g(-x) = 0$.
This only requires $g$ to be an odd function.
Example: $\varphi(z) = e^{z+z^3}$ works for $n=2$ but is not $\exp z$.
\end{solution}

\begin{exercise}
\label{analysis:2011:12}
Use mathematical analysis to show that:
\begin{itemize}
\item[a)] $e$ and $\pi$ are irrational numbers;
\item[b)] $e$ and $\pi$ are also transcendental numbers.
\end{itemize}
\end{exercise}

\begin{solution}
a) For $e$, use $q! e = \text{int} + \sum_{n=q+1}^\infty \frac{q!}{n!}$ which is between 0 and 1.
For $\pi$, use Niven's proof: $f(x) = \frac{x^n(a-bx)^n}{n!}$ and $\int_0^\pi f(x) \sin x dx$.
b) $e$ is transcendental (Hermite, 1873). $\pi$ is transcendental (Lindemann, 1882) using $e^{i\pi} = -1$.
\end{solution}
